\section{Conclusioni, Limitazioni e Sviluppi Futuri}
\label{sec:conclusioni}

Il presente lavoro di tesi ha raggiunto l'obiettivo di tradurre le dinamiche interne dei \textit{Large Language Models} in una rappresentazione visiva e geometrica, implementando un ambiente per l'esplorazione tridimensionale dello spazio latente. L'analisi condotta ha permesso di mappare fisicamente come il contesto alteri le traiettorie di generazione, dimostrando visivamente che le istruzioni ben strutturate agiscono come una mappa cognitiva capace di concentrare la distribuzione dei token verso specifiche sotto-regioni semantiche. 

Pur avendo confermato la possibilità di ispezionare geometricamente la decodifica probabilistica degli LLM, il progetto presenta alcune limitazioni metodologiche e tecniche. In primo luogo, il dataset impiegato per la sperimentazione è stato costruito \textit{ad hoc} e risulta fortemente limitato dal punto di vista numerico; l'assenza di un corpus esteso non ha permesso validazioni statistiche ampie o confronti su casistiche eterogenee. Inoltre, l'analisi condotta non ha previsto uno studio su come la variazione dei parametri del modello o degli algoritmi coinvolti influenzi la topologia dello spazio latente rappresentato. 

Dal punto di vista tecnico, la necessaria riduzione della dimensionalità comporta una fisiologica perdita di granularità semantica. Il sistema di elaborazione e rendering sviluppato, pur rivelandosi perfettamente efficace per lo scopo illustrativo e didattico prefissato, risulta infatti computazionalmente oneroso e limitato nella gestione di vocabolari estesi. Pertanto, i naturali sviluppi futuri di questa ricerca prevedono l'adozione o l'ispirazione a strumenti di livello industriale, come \textit{Nomic Atlas} o \textit{TensorFlow Embedding Projector}, progettati appositamente per visualizzare dataset massivi in modo interattivo e scalabile. Supportata da tali piattaforme avanzate, l'indagine visiva qui iniziata potrebbe espandersi ulteriormente: da un lato integrando le architetture \textit{Retrieval-Augmented Generation} (RAG) per mappare spazialmente l'efficacia del \textit{retrieval} documentale e della \textit{context injection}, dall'altro tracciando i pattern decisionali complessi e i cicli operativi degli agenti.